\documentclass[instructions.tex]{subfiles}
\begin{document}
 
\chapter{Nous devons estimer les afflictions.}
 
Les afflictions sont, dans les desseins de Dieu, des grâces précieuses. Nous devons les estimer, et en remercier sa bonté.
 
\primo Les estimer. Elles sont de grandes grâces pour les justes et pour les pécheurs. Si Dieu destine à une haute prédestination ceux qu'il chérit le plus, il leur fait naître aussi plus d'occasions de souffrir. Oui, ma fille, dit un jour Jésus-Christ à sainte Térèse, plus mon Père envoie d'afflictions à une âme, plus il lui donne de marques de sa tendresse.
 
Les plus grands saints n'ont-ils pas été ceux qui ont le plus souffert en ce monde, et que le monde a plus fait souffrir ? Jésus-Christ, le modèle des prédestinés, le bien-aimé de son Père, n'a-t-il pas été l'Homme de douleur, le plus persécuté et le plus affligé des hommes ? Estimez donc vos afflictions comme des grâces singulières. Dieu vous traite comme ses favoris, en vous rendant plus conformes à son Fils, et plus dignes d'une haute prédestination.
 
Dans les desseins de Dieu, l'adversité est une grâce pour les pécheurs. 1. Dieu les afflige pour les convertir, pour les détacher du monde en leur ôtant ce qui peut les perdre. Il vous ôte la santé, parce que vous en abusez pour l'offenser; il vous ôte vos richesses, parce que vous y êtes trop attaché; il vous a ôté un enfant, parce qu'il était votre idole, et qu'il aurait été le sujet de votre condamnation; il permet qu'on vous décrie, qu'on vous diffame, parce que votre réputation vous remplit d'orgueil et vous aveugle. Dieu vous fait une grâce de vous ôter ce qui peut vous perdre.
 
S'il fallait couper un de vos membres pour vous sauver la vie, vous vous réjouiriez de trouver un homme habile qui vous rendît ce service. Dieu vous rend un service plus grand, lorsque, pour sauver votre âme, il permet qu'il vous arrive quelques disgrâces. N'est-il pas plus avantageux pour vous de voir périr vos biens et souffrir votre corps, que de voir périr votre âme par les délices et par l'abondance ?
 
2. Dieu afflige le pécheur en ce monde pour ne pas le punir en l'autre. Plus l'adversité est grande, plus il devrait estimer la grâce que Dieu lui fait. Loin de se plaindre, il devrait recevoir sa croix avec joie. Il arrive souvent qu'en punition de notre ingratitude et de nos plaintes, Dieu nous ôte, par un effet de sa justice, les croix que son amour nous avait envoyées pour notre salut. S'il nous abandonne à nous-mêmes, s'il nous laisse sans disgrâces, s'il permet que les choses arrivent selon nos désirs, n'est-ce point, hélas! pour notre malheur ? N'avoir point de croix, ah! quelle croix! dit saint Augustin\footnote{Nulla crux, quanta crux! S. Aug.} (1). Un jour viendra que vous n'aurez de consolation que du côté de vos souffrances, et que vous serez bien marris que les choses fussent arrivées autrement.
 
Après tout, qu'est-ce que la souffrance et l'adversité la plus longue en cette vie ? C'est peu de chose quand on regarde l'avenir. Tout ce qu'on souffre ici-bas n'est qu'un moment, quand on le compare à cette éternité de tourmens ou de plaisirs qui doit punir notre lâcheté ou couronner notre patience. Heureux donc ceux qui souffrent avec résignation! plus heureux ceux qui souffrent avec amour de Dieu! mais très-heureux ceux qui souffrent dans le désir de souffrir de plus en plus!
 
\secundo Puisque les adversités sont des faveurs du ciel, nous devons donc en remercier le Seigneur. Voilà une maxime inconnue à la plupart des gens du monde. On les voit aller dans le lieu saint, faire des voyages pour remercier le Seigneur du gain d'un procès, d'une guérison, du succès d'une affaire; mais en voit-on qui lui rendent grâces d'avoir été affligés, ruinés, persécutés ? Ils croient déjà faire beaucoup de ne pas éclater en plaintes contre les auteurs de leurs disgrâces. Cependant point de maxime plus certaine que les croix sont des grâces du ciel. Les saints, quoique aussi foibles et aussi sensibles que nous, en connaissaient le prix, en remerciaient Dieu, les souffraient avec joie; leur patience et leur exemple seront notre condamnation.
 
Saint Paul a été exercé par des contradictions et des épreuves si accablantes, que la vie lui était ennuyeuse et insupportable. Loin de s'en plaindre, je suis, disait-il, rempli de consolation et de joie dans toutes mes tribulations\footnote{Repletus sum consolatione, superabundo gaudio in omni tribulatione nostra. 2. Cor. 7.} (2).
 
Saint Bernard recevait les croix de la main de Dieu avec tant de reconnaissance, qu'il disait : Je serais heureux si j'avais les forces de tous les hommes, afin de porter toutes les croix de l'univers. Sainte Élisabeth, reine de Hongrie, ayant été chassée de son palais indignement traitée par ses sujets, jusqu'à être traînée dans la boue, alla se prosterner devant le Saint-Sacrement, pour remercier Jésus-Christ, et fit chanter le Te Deum en action de grâces.
 
Rien n'était plus fréquent dans la bouche de sainte Térèse que ces paroles : Ou souffrir, ou mourir. La vie lui était insupportable, lorsqu'elle n'avait rien à souffrir pour son Dieu. Saint Jean-de-la-Croix ayant souffert les plus cruels traitemens, le Sauveur lui demanda ce qu'il souhaitait pour récompense de tant de travaux, Seigneur, répondit le saint, je ne vous demande rien que de souffrir de plus en plus pour votre amour.
 
Pourquoi les saints pensaient-ils de la sorte ? C'est qu'ils savaient que plus on souffre avec patience, plus on est semblable à Jésus-Christ, et plus on est aimé de Dieu. C'est pour cette raison que saint Jean Chrysostome, ce grand défenseur des intérêts de Dieu, après avoir souffert les plus dures persécutions, disait : J'aimerais mieux souffrir sur la terre pour la gloire de Jésus-Christ, que de régner avec lui dans le ciel.
 
Les saints aimaient à souffrir, parce que la patience est la preuve de notre amour pour Dieu, et qu'elle ferme l'enfer en nous ouvrant le ciel. Oh! que les maux de cette vie font peu d'impression sur un cœur vivement pénétré de Dieu et des choses de l'autre vie! Tout est facile à supporter quand on aime Dieu. Tout paraît doux, dit saint Bernard, quand il s'agit de mériter une gloire qui ne finira jamais, et d'éviter des supplices qui dureront toujours\footnote{Hæc quàm dulcia meditanti flammas! S. Bern.} (1).
 
\section{Résolutions}
 
\primo Je regarderai toujours les afflictions qui me surviendront ou comme des avertissemens par lesquels Dieu me rappelle à lui, ou comme des moyens qu'il m'offre de me purifier de mes fautes, ou comme des épreuves que je dois subir pour sa gloire, ou comme des occasions de mérite.
 
\secundo Dans toutes ces afflictions, je bénirai Dieu et louerai sa miséricorde. 


\prayer{Mon Dieu, faites de moi ce qu'il vous plaira en cette vie, mais aidez-moi à supporter mes peines, et faites-moi miséricorde pour l'éternité.}
 
\end{document}
